\documentclass[11pt]{article}
\usepackage[letterpaper,margin=1in]{geometry}
\usepackage{amsmath,amssymb,amsthm,mathtools}
\usepackage{microtype,booktabs,xcolor,hyperref,enumitem}
\hypersetup{colorlinks=true,linkcolor=blue!50!black,urlcolor=blue!50!black,citecolor=blue!50!black}
\newtheorem{theorem}{Theorem}[section]
\newtheorem{lemma}[theorem]{Lemma}
\newtheorem{proposition}[theorem]{Proposition}
\newtheorem{corollary}[theorem]{Corollary}
\theoremstyle{definition}
\newtheorem{definition}[theorem]{Definition}
\newtheorem{remark}[theorem]{Remark}
\DeclareMathOperator{\sinc}{sinc}
\DeclareMathOperator{\Tr}{Tr}
\DeclareMathOperator{\ord}{ord}
\DeclareMathOperator{\sig}{sig}
\newcommand{\RR}{\mathbb R}
\newcommand{\dd}{\mathrm d}
\title{What Graph Structure Survives WRT Sewing?\\
\large Hidden connector motifs and the stability boundary of dense cores}
\author{AI Automated Quantum Gravity Research}
\date{}
\begin{document}
\maketitle
\begin{abstract}
We study which graph data survive in the closed, oriented \(\mathrm{SU}(2)\) Witten--Reshetikhin--Turaev invariants of graph manifolds sewn with opposite order-two caps. In square-decorated sewings, an infinite family agrees in its first three logarithmic coefficients but is distinguished by the complete fourth. When the number of connectors grows with level, a centered Neumann kernel governs the first eligible motif contrast. Short-cycle counts do not suffice: a diamond separates connected cubic cores with identical triangle, square and pentagon counts. For a different pair, the actual modulus contrast reverses sign across a scale proportional to \((\log r)^{2/3}\). In rigid dense cores, take the level limit first. After exact vertexwise normalization, any \(o(n)\) maximum-degree deletion is invisible in the limiting signed fluctuation multiplier of a strict fixed template; at linear degree, a loop-permutation contrast survives despite identical local degrees, densities, normalization and exact framing phase. These mechanisms share the same finite oriented sewing but do not interchange the level and size limits.
\end{abstract}
\tableofcontents

\section{Question and regimes}
An invariant of a closed manifold need not record a preferred description of its topology. A sewing presentation still permits a precise question: which changes can the invariant distinguish after fixing elementary data, and which changes are lost? One comparison holds low log coefficients fixed and changes an incidence motif; the other holds local degrees, pair densities, and the exact normalized saddle fixed and changes the composition of permutation masks around a loop. Graph detection here means distinction between specified sewings, not reconstruction of a graph or a canonical JSJ labeling.

In square-decorated sewings integer spin transfer leaves logarithmic charge coordinates. Their correlations resolve shared incidences; the relative order of a pair of invariants can even reverse as graph growth and level compete. In strict dense cores every site instead lies at scale \(\sqrt r\) after the fixed-graph level limit. The subsequent signed continuum saddle remembers a small matrix and tolerates arbitrary inhomogeneous sublinear edge deletion. At linear deletion the matrix changes even when its coarse first-order data are fixed.

We use \(r\) for the shifted level and \(g\) for the number of appended twins. The regimes are (i) fixed square connector, \(r\to\infty\), including every fixed \(g\); (ii) one fixed connector core and \(g=g(r)\to\infty\), with \(g=o(\log r)\) or \(g/\log r\to\alpha>0\); (iii) for a given dense graph \(G_n\), first \(r\to\infty\) and then \(n\to\infty\) in the signed coefficient. The first expansion supplies neither the second nor the third. We distinguish moduli in the square case and signed coefficient and phase information in the dense case.

\section{Closed sewings and exact charge transfer}\label{sec:sewing}
Let \(G\) be finite, connected and simple, and let \(N=|V(G)|\), \(E=|E(G)|\). At a vertex of degree \(d\) take \(\Sigma_{0,d+2}\times S^1\) with section \(z\) oriented as the base boundary and fiber \(h\). Fill the two extra boundaries along \(2z+h,-2z+h\), and glue the remaining ones along backbone edges by \((z,h)\mapsto(h,z)\). These are orientation-reversing boundary identifications. A cycle is a trace in the torus state space, not a division by its dimension. There are no Wilson insertions after filling.

Use \(Z_r(S^2\times S^1)=1\), \(Z_r(S^3)=S_{11}\), \(\tau_r=Z_r/S_{11}\), with
\[
 S_{ab}=\sqrt{2/r}\sin(\pi ab/r),\quad
 \theta_a=e^{\pi i(a^2-1)/(2r)},\quad
 \kappa_r=e^{3\pi i/4-3\pi i/(2r)}.
\]
The core basis is dual orthonormal by the two-point sphere trace. The product-block amplitude is \(\sum_a S_{1a}^{-d}\prod_j S_{ab_j}\) by Verlinde. A filled cap therefore contracts to \((S\Theta^{\pm2}S)_{a1}\); the signs can interchange because a geometric slope \(pz+h\) gives the core ribbon framing \(-p\) in the induced section convention. The pair is unaffected. A seam leaves \(S^3=S\), since the core labels are self-dual.

For the oriented phase, assemble all caps and seam cylinders into a single composition with the disjoint product blocks \cite{BHMV,GM,Turaev}. The test tori use the section Lagrangian \(L_0=\{(q,0)\}\), with \(q,p\) the section and fiber coordinates, \(z\cdot h=1\). Their own kernel is \(L_0\), so individual block, cap and seam tests have zero Maslov correction. Let \(H\) take port sums at each vertex. The product-side actual kernel is \(L_X=\{(H^{\mathsf T}a,p):Hp=0\}\); the base sections have one sum relation per vertex, while fiber copies share one class per vertex. The cap/seam-side kernel is \(L_Y=\{(Tp,p)\}\): a cap has entry \(+2\) or \(-2\), and a seam has the two-by-two block with diagonal zero and off-diagonal \(-1\), from matching cycles across the two ends with opposite boundary signs. Then \(\sig T=0\), \(T\) is invertible, and \(HT^{-1}H^{\mathsf T}=-A_G\). The Maslov form on \((L_X+L_0)\cap L_Y\) is \(p^{\mathsf T}T\tilde p\) on \(\ker H\), using \(x_X=(0,p),x_0=(Tp,0)\). Its signature is \(\sig T-\sig(HT^{-1}H^{\mathsf T})\): the orthogonal complement of \(\ker H\) in the \(T\)-form is \(T^{-1}\operatorname{im}H^{\mathsf T}\), and the residual common radical contributes a hyperbolic space of signature zero. Hence the gluing index is \(\sig A_G\) even when adjacency is singular. In Turaev's unextended composition convention this multiplies the contraction by \(\kappa_r^{+\sig A_G}\): the anomaly unit is \(D/\Delta\), with \(D=1/S_{11}\) and \(\Delta=\sum_a\theta_a^{-1}(S_{1a}/S_{11})^2\). No comparison here requires an absolute Chern--Simons or APS convention.

Put \(K(x)=\sin(\pi x)/(\pi x)\), \(Q(x)=1/K(x)\) on \(0<x<1\), and \(F_d(x)=\cos^2(\pi x/2)Q(x)^d\). For defects \(a=r-\text{color label}\), the paired cap is \(4r^{-1}\cos^2(\pi a/(2r))\) on odd \(a\) and zero on even \(a\). Indeed, for color \(b\) the sine-product expansion of \(S\Theta^2S\) gives, up to a unit factor, \([G(b-1)-G(b+1)]/(2r)\), with \(G(j)=\sum_{\ell\bmod 2r}e^{\pi i(\ell^2+j\ell)/r}\). Translation gives \(G(b+1)=e^{-\pi i b/r}G(b-1)\), and geometric orthogonality in \(|G(j)|^2\) retains only differences \(0,r\), so this norm squared is \(4r\) for \(j+r\) even and zero otherwise. The claimed cap weight follows.
Canceling local charge powers against the edge numerators yields
\begin{equation}\label{eq:wrt}
 \tau_r(M_G)=\frac{\varphi_r(G)}{S_{11}}\,
 2^{2N-E/2}\pi^{-E}r^{3E/2-N} I_r(G),\qquad
 \varphi_r(G)=(-1)^{rE}\kappa_r^{\sig A_G},
\end{equation}
where the real scalar
\[
 I_r(G)=\sum_{\substack{1\le a_i<r\\a_i\text{ odd}}}
       \prod_i F_{\deg_G(i)}(a_i/r)\prod_{ij\in E(G)}K(a_i a_j/r)
\]
is not presumed positive. An upper reflection \(a_i\mapsto s_i=r-a_i\) against an unreflected odd label gives no sign; two adjacent upper reflections give \((-1)^r\). The denominator gains absorb \((s_i/a_i)^{\deg(i)}\), leaving an \(O((s_i/r)^2)\) reflected cap factor times a bounded function for \(s_i\le r/2\).

\subsection{Strict nonbipartite cores}\label{sec:rigidbridge}
For fixed nonbipartite connected \(G\) with weights \(0<w_e<1\), \(\sum_{e\ni i}w_e=1\), we have
\begin{equation}\label{eq:strictbridge}
 (2/\sqrt r)^N I_r(G)\longrightarrow J(G)
 =\int_{(0,\infty)^N}\prod_{ij\in E(G)}K(x_i x_j)\,\dd^Nx .
\end{equation}
At scale \(s_i=\sqrt r\,x_i\), a dyadic box \(x_i\asymp2^{j_i}\) has absolute normalized mass at most \(C2^{\sum_i j_i-\sum_{ij}(j_i+j_j)_+}\), by \(|K(x)|\le\min(1,1/(\pi |x|))\). On a nonempty lattice box the count is \(O(\sqrt r\,2^{j_i})\) at each site. The exponent equals \(\sum_{ij}\{w_{ij}(j_i+j_j)-(j_i+j_j)_+\}\le-c\sum_{ij}|j_i+j_j|\), summable because the edge-sum map is injective on a connected graph with an odd cycle. Upper reflected branches vanish on compact rescaled boxes by the cap gain. The lower odd mesh has spacing \(2/\sqrt r\) and cap factors tend uniformly to one there. This proves absolute convergence and the limit for fixed \(G\), without uniformity in \(N\).

\subsection{Exact reduction of a square}\label{sec:square}
Replace each vertex of a connector \(H\) by an alternating \(K_{2,2}\) and join the designated vertex of each square along connector edges, obtaining \(\square(H)\). For odd \(a,c<r\), set \(s_a=\min(a,r-a)\). The degree-two sum
\begin{equation}\label{eq:two}
 \sum_{b<r,\ b\text{ odd}}F_2(b/r)K(ab/r)K(cb/r)
 =\frac{r}{4ac}\min(s_a,s_c)
\end{equation}
follows by extending to the full odd \(2r\)-grid and dividing by two; \(b=r\), if present, gives zero. The summand before \(1/(4ac)\) is \(\sin(\pi ab/r)\sin(\pi cb/r)/\sin^2(\pi b/(2r))\). The cosine difference on the full \(M\)-grid is \(2j(M-j)\) by geometric orthogonality; subtracting the even terms gives \(2r\min(j,2r-j)\) on the odd grid. Product-to-sum gives the result. Applying this to the two opposite sites in each block and summing the nonconnector site gives
\begin{align}
 W_r(a)&=\sum_{c<r,\ c\ {\rm odd}}F_2(c/r)c^{-2}\min(s_a,s_c)^2,\quad
 p_{r,d}(a)=a^{-2}F_{2+d}(a/r)W_r(a),\notag\\
 P_H(r)&=\sum_{\boldsymbol a<r,\ {\rm odd}}\prod_i p_{r,d_i}(a_i)
        \prod_{ij\in E(H)}K(a_i a_j/r),\qquad
 I_r(\square(H))=(r/4)^{2m}P_H(r),\label{eq:square}
\end{align}
where \(m=|V(H)|\), \(d_i=\deg_H(i)\). For \(q=|E(H)|\), \(\chi=q+1\), \(L=\log r\), and \(h_0=1/2\), define
\[
 A_{m,q}(r)=2^{m-\chi/2}\pi^{-4m-\chi}r^{4m+3\chi/2}.
\]
Then exactly \(|\tau_r(M_{\square(H)})|/A_{m,q}(r)=\pi h_0^{-m}|P_H(r)|/[r\sin(\pi/r)]\), valid at either parity even in a joint limit.

\section{Fixed connectors: topology beyond three coefficients}\label{sec:fixed}
Let \(H_0,G_0\) be bipartite \(4+4\) graphs whose A-to-B row masks, with bit \(j\) referring to \(B_j\), are
\[
 H_0=(3,7,11,12),\qquad G_0=(3,5,14,14).
\]
Append \(g\ge1\) distinct B sites, each adjacent to all four A sites, obtaining \(H_g,G_g\). These are connected, leafless and nonisomorphic: the two original degree-three B sites have identical neighbor sets in one graph and different sets in the other, and the two bipartition sizes \(4,4+g\) are unequal. Denote \(\lvert\tau_r(M_{\square(H)})\rvert\) by \(W_r[H]\); this is distinct from the function \(W(s)\) below.

\begin{theorem}[Fixed fourth-log distinction]\label{thm:fixed}
For every fixed \(g\ge1\), the normalized moduli \(W_r[H_g]/A_{m,q}(r)\) and \(W_r[G_g]/A_{m,q}(r)\), with \(m=8+g,\ q=10+4g\), have expansions through \(L^{m-3}\). The first three coefficients coincide. The full fourth-coefficient difference is
\[
 C_4(H_g)-C_4(G_g)=\frac{g\,\mathcal R}{(g+2)(g+3)(g+4)(g+5)}>0,
 \qquad \mathcal R=2\int_0^\infty W(s)\bigl(W(s)-s\bigr)\dd s,
\]
where \(W(s)=\int_0^\infty K(x)\bigl(K(e^{-s}x)-1\bigr)\dd x/x\).
The result holds as \(r\to\infty\) over all integer shifted levels separately for each \(g\).
\end{theorem}

\begin{proof}
We first compare the entire cap and lattice correction, then subtract the continuous logarithmic products before computing their coefficients.

\paragraph{Finite weights versus the continuous reference.}
For odd \(a\ge1\) define
\[
 p(a)=\sum_{c\ge1,\ c\ {\rm odd}}\max(a,c)^{-2},\quad
 c_p=\tfrac12(\gamma+\log2)+\tfrac14,\quad
 \mu_p=\sum_{a\ {\rm odd}}p(a)\delta_{\log a},\quad
 \mu_0=h_0\,\dd t+c_p\delta_0 .
\]
Here \(\gamma\) is Euler's constant. Counting the odd pairs with maximum \(a\) gives
\(\sum_{a\le X,\ {\rm odd}}(p(a)-1/a)=\#\{a\le X:a\text{ odd}\}
\sum_{c>X,\ {\rm odd}}c^{-2}=1/4+O(1/X)\).
Thus \(D(t)=(\mu_p-\mu_0)([0,t])=O(e^{-t})\), and both measures have uniformly bounded unit-interval masses. For fixed connector degree \(d\), \eqref{eq:square} yields \(p_{r,d}(a)=p(a)+O(r^{-1})\) on the lower half: \(W_r(a)\) differs from its infinite \(F_2=1\) counterpart by \(O(a^2/r)\). After all incident reflections the upper weight is \(O(s_a/r^2)\). These bounds apply in any mixture of replaced and unreplaced weights.

For a leafless fixed connector define
\[
 V_H=\operatorname{vol}\{0\le x_i\le1:\ x_i+x_j\le1\ (ij\in E(H))\},
 \qquad \mathcal D_H=\sum_{\deg_H(v)=2}V_{H-N[v]} .
\]
Volumes include isolated surviving sites and assign volume one to the empty set. Set \(\Phi(u)=K(e^{-u})\). The difference between \(h_0^{-m}P_H(r)\) and the reference integral on \([0,L]^m\) with measure \((\mu_0/h_0)^{\otimes m}\) and edge weights \(\Phi(L-t_i-t_j)\) is
\begin{equation}\label{eq:micro}
 \Gamma\,\mathcal D_H L^{m-3}+o(L^{m-3})
\end{equation}
for a universal \(\Gamma\). For a degree-\(d\) cap replacement, retaining the \(d\) adjacent sine bounds costs \(O((1+L)^{m-d-1})\): their neighboring harmonic sums are \(O(1+\log(r/s_a))\), and \(r^{-1}\sum_{a\le r/2}(1+\log(r/a))^d=O(1)\). Thus degree at least three is negligible. For degree two, condition on its neighbors \(b,c\). The conditional difference tends, at fixed odd \(b,c\), to
\[
 h_0\int_0^1\frac{M(x)-1}{x}K(bx)K(cx)\dd x,\quad
 M(x)=F_4(x)\,j(\min(x,1-x))/x,\quad
 j(u)=\pi\int_0^u v\cot(\pi v)\dd v .
\]
Indeed \(W_r(rx)/r\to j(\min(x,1-x))\) by the odd-mesh \(\min\)-square integral, whose derivative for the lower half is \(\pi x\cot(\pi x)\). Upper reflection of the conditioned site absorbs both incident denominators, giving \(O(s_a/r^2)\) instead of \(O(1/s_a)\). Splitting the conditioned sum at \(r/\max(b,c)\) gives \(C(1+\log\max(b,c))/\max(b,c)\) after reflection, summable against \(1/(bc)\). Upper branches of the two neighbors vanish pointwise after their incident reflections. With these charges fixed the other normalized coordinates give \(V_{H-N[v]}\): constraints to a fixed charge affect only the top boundary of the exponent cube, and any edge between the fixed charges tends to one.

Next telescope \(\mu_p\) to \(\mu_0\) and integrate \(D\) by parts. If the differentiated kernel's high neighbor at \(t\) has degree at least three, retaining two other edges gives \(O((1+L-t)^2)\) after integrating their other coordinates. The product with \(\min(1,e^{2(s+t-L)})|D(s)|\) leaves only \(O(L^{m-4})\). At a degree-two high neighbor write its coordinate \(L-y\), and the opposite low coordinate \(z\); the local factor is \(D(s)\Phi'(y-s)\Phi(y-z)\). The high-site mesh tends to \(h_0\dd y\), the same deletion volume appears, and the \(z,y\) integrals are dominated against \(O((1+s)^2)|D(s)|\). The two derivative origins see \(\mu_p,\mu_0\) once each at the opposite low coordinate, independent of the telescoping order. With \(B_p(y)=\int_0^\infty\Phi(y-z)\mu_p(\dd z)\) and \(B_0\) defined analogously, their sum uses \(\int(B_p^2-B_0^2)\dd y\), since \(B_p-B_0=\int D(s)\Phi'(y-s)\dd s\). Combining the mesh and cap terms in the \(h_0^{-m}\) normalization gives
\[
 \Gamma=4\int_0^\infty \bigl(M(e^{-y})B_p(y)^2-B_0(y)^2\bigr)\dd y.
\]
This is integrable by \(M(e^{-y})=1+O(e^{-2y})\), \(B_p-B_0=O(e^{-y})\) at infinity, and the upper cap cancellation at zero. The leading \(V_H L^m>0\) allows removal of the modulus at all displayed orders; \(\pi/[r\sin(\pi/r)]\) affects none of them.

The only degree-two sites of either graft are old B sites. Their deletion invariants agree: reverse B coordinates to \(x_A\le y_B\), where each linear extension contributes volume \(1/(m-3)!\) after deletion. In an old order let \(j\) count sites after the final surviving A; insert the \(g\) labeled twins in \((g+j)!/j!\) ways. By appending unused vertices only when their predecessors have appeared, the common old deletion histogram is \((0,0,4,24,0)\) for \(j=0,\ldots,4\), checked also by the supplied integer certificate. Thus \eqref{eq:micro} cancels.

\paragraph{Ordered-gap cancellation.}
Reverse B coordinates and sort the \(m-a\) free values after choosing \(a\) atoms. A pins lie at 0 and B pins at \(L\). Each edge spans an interval of consecutive gaps. Forward edges give \(1+f_+\), backward edges \(f_-\), where \(f_+(u)=\Phi(u)-1\), \(f_-(u)=\Phi(-u)\) for \(u\ge0\), evaluated at the sum of that edge's gaps. An edge between pins gives \(1+O(e^{-2L})\). If the chosen defect intervals use \(b'\) gaps, integrate out the others at unit Jacobian to get
\((L-\sum z)_+^{m-a-b'}/(m-a-b')!\) times the product of the \(f\)'s in the used gaps. A term using every gap is exponentially small; all products decay exponentially in the sum of used gaps. Consequently coefficients through order \(m-3\) require only \(a+b'\le3\).

For reproducibility, enumerate free orders and pin sets, include every backward defect, and insert optional forward defects while their union has at most \(3-a\) gaps. Store the multiset of signed nonempty gap masks, canonically minimizing over permutations of the selected gap labels. Identical lists have identical integrals. Only \(d=0,1,2,3\) prototype twins are needed: each touched twin is pinned or in a nontrivial gap-contracted block, costs a distinct pin or gap, and every nontrivial free block contains an A. An untouched twin must be after the last A-block. If \(j\) old B singletons follow that block, insertion has multiplicity
\(\binom gd(g-d+j)!/j!=g!\binom{g-d+j}j/d!\).
If the last A is absorbed into the fixed right endpoint no untouched twin fits; at most two can be touched, and \(g=1,2\) are checked directly. The finite signed count of the stated masks is
\[
(d,j)=(0,1): -2\mathcal S,\qquad (d,j)=(0,2): 2\mathcal S,
\qquad \mathcal S=\operatorname{sig}_{\rm gaps}(C_4-2P_4+M_2),
\]
in the no-atom signature, with all other sectors zero. Here \(M_2\) denotes two disjoint edges, and \(\mathcal S\) uses exactly three gaps. The bundled C++/Python replay checks the entire integer signature, including \(g=1,2\); this is a finite exact count, not a numerical integral. The insertion multiplier yields \(g!g(g+1)\mathcal S\), so only the fourth continuous coefficient differs, with denominator \((g+5)!\).

To evaluate its integral, take ordered A coordinates \(p,p+s\) and put \(q'=L-p-s\). The B transfers for both neighbors and for a single high or low neighbor are
\(H_{\mathrm{tr}}=\int_{-p-s}^{q'}\Phi(u)\Phi(u+s)\dd u\),
\(J_n=\int_{-p-s}^{q'}\Phi(u)\dd u\), and
\(J_f=\int_{-p}^{q'+s}\Phi(u)\dd u\).
Including both A orders, the contrast has integrand
\(2(H_{\mathrm{tr}}-J_n)(H_{\mathrm{tr}}-J_f)\).
The exponential tails of \(\Phi-\mathbf1_{u\ge0}\) give differences
\(W(s)+O(e^{-p-s}+e^{-2q'-2s})\) and \(W(s)-s+O(e^{-2q'}+e^{-p})\).
Thus the integral is \(L\mathcal R+O(1)\). Finally the integral of \(K(x)(1-\cos(\pi z x))\dd x/x\) for \(0\le z\le1\) equals
\(\tfrac12[(1+z)\log(1+z)+(1-z)\log(1-z)]\), by inserting damping, differentiating, integrating the sine-product Frullani formula, and removing damping. Averaging with \(z=e^{-s}v\) for \(v\in[0,1]\) shows \(W(s)<0\) and \(W(s)=O(e^{-2s})\). This gives \(\mathcal R>0\) and completes the coefficient calculation.
\end{proof}

\section{Growing connectors: when leading data forget incidence}\label{sec:growing}
Let a fixed bipartite core have \(n\) A labels and \(b\) old B sites with nonempty neighborhoods \(S\subseteq[n]\), and append \(g\) universal B sites. The number \(q\) of connector edges may now depend on \(g\). Define the character \(J_a(x)=\sin(\pi a x)/(a\sin(\pi x))\) and the B transfer
\[
 T_{r,S}(\boldsymbol a)=\sum_{c<r,\ c\ {\rm odd}} p_{r,0}(c)\prod_{i\in S}J_{a_i}(c/r),
 \qquad B_r(\boldsymbol a)=\Bigl(\prod_i Q(a_i/r)\Bigr)T_{r,[n]}(\boldsymbol a).
\]
Conditioning on A in \eqref{eq:square} gives the exact product of each old \(T_{r,S}\), the common \(B_r^g\), and the individual \(p_{r,d_i}\), with \(d_i\) the old A degree. Reflect upper A labels before multiplying these factors. At every old B label \(c\) odd, \(K((r-s)c/r)=(s/(r-s))K(sc/r)\) and \(Q((r-s)/r)s/(r-s)=Q(s/r)\). Thus \(B_r\) is unchanged as a function of the reflected labels, and the reflected A measure has lower contribution \(p_{r,d_i}(s)\) and upper contribution \(p_{r,d_i}(r-s)(s/(r-s))^{d_i}\), where the original labels are odd and the reflection midpoint is counted only once. The upper weight is \(O(s/r^2)\) as in Section~\ref{sec:fixed}. In particular the exact high degree of the newly added sites does not enter the sitewise constants.

For \(\boldsymbol s\) reflected, write \(M_S=\max_{i\in S}s_i,\ M=M_{[n]}\). Uniformly,
\begin{equation}\label{eq:bound}
 T_{r,S}(\boldsymbol s)=h_0(L-\log M_S)+O(1),\qquad
 B_r(\boldsymbol s)=h_0(L-\log M)+O(1).
\end{equation}
For \(c\le r/2\) the exact weight is \(p(c)+O(1/r)\); use \(|J_a|\le1\), \(|1-J_a(x)|\le C\min(1,(ax)^2)\) on the lower half, and \(|J_a(x)|\le C/(ax)\) there. Split at \(c=r/M_S\): replacing the kernel by one below the cut and bounding it with the maximal charge above each cost \(O(1)\) in the harmonic mass. The upper \(c\) contribution is bounded by \(p_{r,0}(c)=O((r-c)/r^2)\), with the characters bounded by one. The ratio \(Q(s_i/r)=1+O((M/r)^2)\) costs at most \(O(1)\) times the logarithmic estimate. This proves \eqref{eq:bound}, with no positivity assertion at finite \(r\).

For fixed odd \(\boldsymbol a\), define
\[
 \rho(x)=\frac{F_2(x)j(\min(x,1-x))}{x^2},\qquad
 \mathcal C_S(\boldsymbol a)=c_p+h_0\int_0^1
       \Bigl[\rho(x)\prod_{i\in S}J_{a_i}(x)-\frac1x\Bigr]\dd x .
\]
The integral converges since \(\rho(x)=1/x+O(x)\) near zero and \(\rho(x)=O(1-x)\) at the other end. Subtracting the harmonic mass and using the odd Riemann sum and the endpoint bounds of Section~\ref{sec:fixed} gives \(T_{r,S}-h_0L\to\mathcal C_S\). For this fixed-label limit \(\prod Q(a_i/r)=1+O(r^{-2})\) as well.

\begin{proposition}[Two distinct charge localizations]\label{prop:thermal}
If \(g/L\to\alpha>0\), the signed reduced scalar divided by \((h_0L)^{g+b}\) tends to the positive Gibbs mass
\[
 Z_n(\alpha)=\sum_{\boldsymbol a\ {\rm odd}\ \ge1}
       \prod_i p(a_i)\exp\bigl(\alpha\mathcal C_{[n]}(\boldsymbol a)/h_0\bigr).
\]
In the separate regime \(g\to\infty,\ g=o(L)\), put \(T=L/g\). The same normalized scalar divided further by \(T^n\) tends to \(h_0^n n!\), with normalized limiting A-charge density \(e^{-\max_i u_i}/n!\) for \(u_i=\log(s_i)/T\). In either regime the leading normalized invariant has no dependence on the internal incidence of the fixed core. When comparing raw values the prefactor \(A_{m,q}(r)\) must be restored.
\end{proposition}
\begin{proof}
At finite \(\alpha\) the common factor satisfies
\(\lvert B_r/(h_0L)\rvert^g\le C_\alpha M^{-\alpha/2}\), using \eqref{eq:bound} and \(1+x\le e^x\). Together with the A measure bound \(C/s_i\) this is summable, even with any fixed power of \(1+\log M\), by grouping according to \(M\). The upper A labels vanish pointwise, and the fixed lower ones give the stated Gibbs mass with a positive all-one term. In the mesoscopic regime the common factor goes to \(e^{-\max u_i}\) on compact scaled sets and has absolute bound \(Ce^{-\max u_i}\) globally, again by \eqref{eq:bound}. Each old transfer divided by \(h_0L\) tends to one on compacts. The measures \(\mu_p/T\) in the scaled log coordinate converge to \(h_0\dd u\) and have uniformly bounded masses on scaled unit intervals; the upper contribution and lower errors vanish on each compact with the same tail domination. Layer cake in the radial maximum gives \(\int e^{-\max u}\dd^nu=\int_0^\infty t^n e^{-t}\dd t=n!\). Both proofs apply to bounded tests and show that the signed real scalars are eventually positive.
\end{proof}

\subsection{A finite-subtraction comparison and a positive motif law}
A useful transparent family begins with a simple \(d\)-regular graph \(J\) of \(n\) vertices. Its vertices are the A sites; each old edge contributes one B site \(S=\{i,j\}\). One may add a common old universal B site in both cores to ensure connectivity. Then append \(g\) universal B twins. Write \(H_g,G_g\) for the two resulting connectors from \(J_H,J_G\) with the same \(n,d\). The exact reflected A product is invariant under any permutation of labels, since the original A degrees agree sitewise. In the finite product expand each old \(T_{r,S}\) as \(h_0L+X_{r,S}\). A term selects \emph{distinct} old B sites, leaving shared A variables shared. An old universal site has the same symmetric factor in each sum. Every lower sector cancels at finite \(r,g\) if the counts of the selected edge shapes in that sector agree, including any shared vertices and isolated padding. This prevents errors from subtracting two separately approximated leading terms.

\begin{theorem}[First eligible mesoscopic comparison]\label{thm:motif}
Suppose the selected-edge shape counts agree at every order below \(\ell\). Set
\[
 Z_J(z)=\int_{[0,1]^n}\prod_{uv\in E(J)}[1-z\max(x_u,x_v)]\,\dd^n x,\qquad
 C(x,y)=\tfrac12(x^2+y^2)-\max(x,y)+\tfrac13 .
\]
For \(g\to\infty,\ g=o(L)\),
\begin{equation}\label{eq:motif}
 g^\ell\frac{W_r[H_g]-W_r[G_g]}{W_r[G_g]}
 \longrightarrow \frac{(n+\ell)!}{n!}[z^\ell](Z_{J_H}-Z_{J_G}).
\end{equation}
The coefficient on the right is a sum over \(\ell\)-edge motifs with no degree-one vertices:
\[
 [z^\ell](Z_{J_H}-Z_{J_G})
 =\sum_F (N_{J_H}(F)-N_{J_G}(F))\,\int\prod_{uv\in E(F)}C(x_u,x_v)\,\dd^{|V(F)|}x .
\]
Every weight in this sum is strictly positive. An exact signed cancellation can still occur between different motifs. For \(g/L\to\alpha>0\), the analogous first eligible limit scales by \(L^\ell\) instead of \(g^\ell\) and equals \(h_0^{-\ell}\) times the Gibbs expectation of the degree-\(\ell\) difference of \(\prod_{\rm old\ }(1+z\mathcal C_S)\).
\end{theorem}
\begin{proof}
Under \eqref{eq:bound}, on the scaled log domain \(X_{r,S}/(h_0T)\to-\max_{i\in S}u_i\). Its polynomial growth in the maximum is absolutely dominated by the common exponential bound. Degree \(j\) sectors, after normalizing by the leading mass, are \(O(g^{-j})\), so the exactly canceled sectors below \(\ell\) cannot reappear as measurement or approximation errors. The remaining homogeneous polynomial at degree \(\ell\) gets radial factorial \((n+\ell)!\) from the same layer cake as in Proposition~\ref{prop:thermal}. At \(g/L\to\alpha>0\) one instead has \(X_{r,S}\to\mathcal C_S\) on each fixed-label configuration, and the summable Gibbs envelope tolerates its fixed powers; this yields the last assertion. In both cases the physical prefactors match and the reduced scalars are eventually positive.

For the centering, let \(U_z(x,y)=1-z\max(x,y)\). There is a positive continuous \(a_z\), analytic near \(z=0\), satisfying
\[
 \int U_z(x,y)a_z(y)^{d-1}\dd y=a_z(x)N_z,\qquad N_z=\int a_z^d,\qquad a_0=1 .
\]
The derivative of right minus left at the constant \(1\) is \(f\mapsto f+\int f\), a bounded isomorphism on \(C[0,1]\), so the implicit function theorem applies. With \(\mu_z=a_z^d\dd x/N_z\) the kernel \(B_z=U_z/(a_z\otimes a_z)-1\) has exactly zero mean in each variable. By regularity \(Z_J=N_z^n\int\prod_{uv}(1+B_z(x_u,x_v))\prod_i\mu_z(\dd x_i)\); a chosen subgraph with a leaf integrates to zero. Differentiating at zero gives \(B_z=zC+O(z^2)\) uniformly. Lower shape agreement gives the displayed coefficient formula. Agreement only of lower leafless motif counts suffices for this hard polynomial, but has not been used in the actual finite WRT cancellation.

The kernel \(C\) is the mean-zero Neumann Green kernel, as follows from \(-\partial_x^2 C=\delta_y-1\) and zero end derivatives and zero mean. Its absolutely uniform expansion is \(\sum_{m\ge1}(\pi m)^{-2}e_m(x)e_m(y)\), \(e_m(x)=\sqrt2\cos(m\pi x)\). At a vertex every integral of incident cosine products is a nonnegative multiple of the number of sign choices with signed frequencies summing to zero. A strictly positive term exists for a connected leafless graph: choose a signed incidence matrix with a full-support rational kernel vector. When there are no bridges, use ordinary oriented incidence since removing any edge leaves rank unchanged. With bridges, remove them and consider the resulting block tree. Each leaf block contains a cycle (otherwise a singleton leaf block would have original degree one). Mark one edge of a cycle in each leaf block negative, others positive, using same-sign incidence at the two ends of a negative edge. Each side of any bridge has an unbalanced cycle and any nonbridge deletion leaves some unbalanced cycle intact. The connected signed-incidence rank is full in the presence of such a cycle, by propagating any left-null vector along the edges. Hence any one column removal leaves its rank unchanged, so there is a rational null vector nonzero on every column (avoid the finite union of proper coordinate hyperplanes). Clear denominators and use the magnitudes as integer frequencies, taking the local signs from the signed incidence relation. Disconnected motifs factor over components. This proves strict positivity.
\end{proof}

For example, if two \(d\)-regular cores have girth at least \(\ell\), their lower shape counts coincide: any source homomorphism with fewer edges lifts into the universal \(d\)-regular tree once its root is lifted, because all fundamental cycles map to walks too short to wind. Inclusion--exclusion on coincident labels turns these hom counts into injective ones and tolerates quotient loops (zero) and repeated edges (one constraint). At \(\ell\) the only leafless possibility is an \(\ell\)-cycle, whose weight is \(\Tr(C^\ell)=\zeta(2\ell)/\pi^{2\ell}\). Thus \(C_\ell\sqcup C_\ell\) versus \(C_{2\ell}\) gives arbitrarily deep first separating orders in the mesoscopic regime, with the common old universal B site making both connectors connected.

With shorter cycles present, counts of cycles alone can fail at the first eligible order. In a regular core triangle and four-cycle equality implies all selected-edge shape counts match through four: remove source leaves (factor \(d\)) and isolated vertices (factor \(n\)); the only nonempty leafless cores remaining are triangles and squares, and the noninjective part of the square count is governed by smaller quotients. Then use the same inclusion--exclusion. At order five the only possibilities are the pentagon and the diamond \(D=K_4-e\); their weights are \(1/93555\) and \(1/467775\). The latter follows by integrating \(C(x,y)(C^2(x,y))^2\dd x\,\dd y\), where \(C^2(x,y)=\int C(x,u)C(u,y)\dd u\); splitting the inner polynomial at \(x,y\) and doubling the ordered triangle uses \(\int_{0\le x\le y\le1}x^iy^j\dd x\dd y=1/[(i+1)(i+j+2)]\). The signed hard coefficient is therefore
\begin{equation}\label{eq:diamond}
 \frac{5\Delta c_5+\Delta N(D)}{467775}.
\end{equation}
Here is a matched cubic pair \(D_H,D_G\) on vertex labels \(0,\ldots,11\):
\[
\begin{array}{c|l}
D_H&(0,1),(0,4),(0,8),(1,4),(1,9),(2,5),(2,7),(2,11),(3,6),\\
 &(3,7),(3,11),(4,9),(5,8),(5,11),(6,7),(6,10),(8,10),(9,10)\\
D_G&(0,4),(0,5),(0,9),(1,2),(1,9),(1,11),(2,6),(2,9),(3,4),\\
 &(3,7),(3,8),(4,5),(5,7),(6,10),(6,11),(7,8),(8,10),(10,11).
\end{array}
\]
Both have \((c_3,c_4,c_5)=(4,2,4)\), while summing \(\binom{|N(u)\cap N(v)|}{2}\) over edges gives one diamond in \(D_H\) and none in \(D_G\). Thus the full mesoscopic modulus contrast is positive at \(g^{-5}\), with coefficient \(17!/(12!\,467775)\). The exact edge lists and the polynomial integrations are also checked by the replay.

\subsection{A sign reversal from competing hard and soft motifs}\label{sec:reversal}
Construct two further cubic graphs on 12 labels:
\[
\begin{array}{c|l}
C_H&(0,3),(0,4),(0,6),(1,4),(1,10),(1,11),(2,8),(2,9),(2,11),\\
 &(3,5),(3,10),(4,7),(5,9),(5,11),(6,7),(6,8),(7,9),(8,10)\\
C_G&(0,3),(0,5),(0,10),(1,3),(1,4),(1,5),(2,6),(2,8),(2,10),\\
 &(3,7),(4,9),(4,10),(5,11),(6,7),(6,11),(7,8),(8,9),(9,11).
\end{array}
\]
They have no triangles or diamonds; their four-cycle counts are both two, but their five-cycle counts are five and six. Let \(J_H\) consist of five copies of \(D_H\) and one \(C_H\), and let \(J_G\) be the corresponding G copies. These edge graphs have \(n=72\), \(b=108\). Add one common old universal B site so the bipartite cores are connected, and then graft \(g\) twins as before.

\begin{theorem}[Actual modulus reversal]\label{thm:reversal}
For the specified cores, as \(g\to\infty,\ g=o(L)\) their actual closed WRT moduli obey
\[
 \frac{W_r[H_g]-W_r[G_g]}{W_r[G_g]}
 =\mathcal A g^{-6}+\mathcal B g^{-5}T^{-2}
    +o(g^{-6}+g^{-5}T^{-2}),\qquad T=L/g,
\]
where
\[
 \mathcal A=\frac{2449558508}{945}>0,\qquad
 \mathcal B=\frac{27010}{63}\left(\frac{\pi^2}{24}-\frac12\right)<0 .
\]
The sign is positive when \(g/L^{2/3}\to0\) and negative when \(g/L^{2/3}\to\infty\) within this regime. At finite limiting ratio \(\lambda\), the \(g^6\)-scaled gap has limit \(\mathcal A+\mathcal B\lambda^3\), with no sign asserted when that number is zero.
\end{theorem}

\begin{proof}
All selected-edge shapes through four match. By \eqref{eq:diamond} and \(\Delta c_5=-1,\ \Delta N(D)=5\), the hard coefficient of degree five also cancels. For the hard sixth coefficient an exact useful recurrence is obtained by integrating on \([0,t]^S\), retaining edges within \(S\), and splitting by the maximal site \(v\). If the coefficient of \(z^q\) is \(c(S,q)t^{|S|+q}\), then
\[
 c(S,q)=\frac1{|S|+q}\sum_{v\in S}\sum_{j=0}^{\min(q,d_S(v))}
             (-1)^j\binom{d_S(v)}j c(S\setminus v,q-j),\qquad c(\varnothing,0)=1 .
\]
Use \(c(\varnothing,q)=0\) if \(q>0\); \(d_S(v)\) counts incident edges in \(S\). For the two component contrasts the degree-six coefficients are \(-3911/212837625\) and \(53/500500\), the degree-five coefficients \(1/467775\) and \(-1/93555\), and the common degree-one coefficient of each polynomial is \(-12\). Multiplying over the six components gives \([z^6]\Delta Z_J=11933/851350500\), since all cross contributions from the degree-five differences have the same multiplier and cancel. Layer cake gives \(\mathcal A=(78)!/72!\) times this rational coefficient. Selecting the common old universal B at sixth order leaves a degree-five angular difference, already zero, times the common radial maximum. Orders seven and above are \(O(g^{-7})\) after leading normalization.

To obtain the competing fifth coefficient, retain the exact radial weight until all hard angular cancellations have occurred. Put \(t_i=\log s_i\), \(t=\max t_i\), and cut at \(t_*=L-\sqrt L\). Outside this body \eqref{eq:bound} bounds the normalized common power by \((2/\sqrt L)^g\) for large \(r\); the other masses cost only fixed powers of \(L\). The error is smaller than any prescribed negative power of \(L\) for every \(g\to\infty\). In the body upper A weights and the replacement of the lower A weights by \(p\) cost relatively \(O(e^{-\sqrt L})\), since \(s/r\le e^{-\sqrt L}\). Use a common shifted old reference \(h_0L+c_\infty\), with \(c_\infty=c_p+h_0\int_{\mathbb R}(\Phi(u)-\mathbf1_{u\ge0})\dd u\); lower cancellations are unchanged. On this body the needed transfers are
\begin{align}
 T_{r,\{i,j\}}-h_0 L-c_\infty&=-h_0\max(t_i,t_j)+h_0W(|t_i-t_j|)+R_{ij},\notag\\
 B_r-h_0 L-c_\infty&=-h_0t+h_0\sum_{i\ne i_{\max}}W(t-t_i)+R_* .\label{eq:tieexp}
\end{align}
Here \(R_{ij}=O((1+t_{\{i,j\}})e^{-t_{\{i,j\}}}+e^{-\sqrt L})\) and \(R_*\) has the analogous one-site error plus \(O(Le^{-2\sqrt L})\) and a sum of \(O(e^{-\max(t-t_i,t-t_j)})\) for distinct nonmaximal \(i,j\).
For clarity, the estimates follow by writing the character product around its maximum \(A\) as \(J_A\prod(1+(J_B-1))\). For the singleton, replace \(p_{r,0}\) by \(p\) in the lower half, \(J_A(x)\) by \(K(Ax)\), and extend the \(p\)-sum; the respective errors are bounded by \(C(1+\log A)/A+C A/r\) using the small-argument and \(1/(Ax)\) bounds in \eqref{eq:bound} (the reflected upper half has \(O((r-c)/r^2)\) weight). Integration by parts in the cumulative discrepancy \(D\) costs \(O(A/r)\), by
\(\left|\partial_{\log c}K(Ac/r)\right|\le C\min((Ac/r)^2,1)\);
the atom and continuum lower limit give the constant \(c_\infty\). With one difference \(J_B-1\) the same bounds hold: the discrepancy derivative for \(K(Ac/r)(K(Bc/r)-1)\) obeys the same estimate uniformly for \(B\le A\), and the continuum term after rescaling is \(h_0W(\log(A/B))\). For products containing at least two differences retain the smallest charge \(B\) and the maximal \(A\) in the bound: split \(c<r/A,\ r/A\le c<r/B,\ c\ge r/B\), to get \(O(B/A)\). This localizes at least two gaps. The immediate \(Q\) correction is \(O(Le^{-2\sqrt L})\). This proves \eqref{eq:tieexp}. The explicit cosine integral for \(W\) in Theorem~\ref{thm:fixed} also yields
\[
 w_0=\int_0^\infty W(s)\dd s
   =-\sum_{j\ge1}\frac1{(2j)^2((2j)^2-1)}=\frac{\pi^2}{24}-\frac12 ,
\]
and \(W,W'\) decay exponentially (the derivative series is summable also at \(s=0\)).

In an unnormalized fifth moment, a single \(W\) loses one free log by tying a pair, and replaces one growing factor (or costs \(T^{-1}\) when differentiating the common power). Its possible scale is \(T^{n+3}\), against \(T^{n+5}\) for the hard fifth term. With bounded unit log masses and common bound \(Ce^{-t/T}\), two soft factors lose at least one log while replacing two powers, hence give \(O(T^{n+2})\). The remainders in \eqref{eq:tieexp} either localize at least two logs or differences, or are negligible on the body; they give the same bound. The first two derivatives of the normalized common power with respect to the bounded base shift satisfy \(CT^{-j}e^{-t/T}\), \(j=1,2\), since the exponent is \(g\) and the base before normalization is \(L+O(1)-t\) with \(t\le t_*\).

For the pure hard lattice comparison replace \(\mu_p\) by \(h_0\dd t+c_p\delta_0\) one coordinate at a time. A derivative of a hard edge maximum with respect to \(t_i\) vanishes unless the other endpoint is \(\le t_i\); thus the discrepancy \(D(t_i)\) exponentially pins \emph{two} coordinates and the derivative removes one power of \(T\). A radial derivative anchors all coordinates below \(t_i\) and costs \(T^{-1}\). If some anchored coordinates have already been pinned their log powers were already lost. Therefore the cumulative replacement error is \(O(T^{n+2})\). The continuous and the one- and two-atom contributions are zero: for every \(J\) the hard cube restricted to \([s,1]^n\) is
\[
 (1-s)^n(1-zs)^bZ_J(z(1-s)/(1-zs)).
\]
At \(s=0\), the first derivative is minus the sum of one-pin face integrals. The second derivative is twice the sum of two-pin integrals: the normal integrand derivative on a lower face vanishes a.e., since each neighboring argument is larger. Equality of the hard polynomials through order five thus kills these summed face differences, and by homogeneity their cancellations hold at each radius, including with the exact finite-\(g\) radial weight. Three or more pins are \(O(T^{n+2})\).

For a linear tie, replacing all measures by \(h_0\dd t\) costs a further log: an atom or discrepancy anchors a changed coordinate as well as the tied gap, or if that coordinate belongs to the tie it anchors both partners. Integration by parts uses the exponential bounds on \(D,W,W'\); any ordering-boundary point, including a lattice equality, has the same loss, since a jump is evaluated with the anchored discrepancy or with lattice atom size \(O(e^{-t_i})\). On each of the two tie orientations substitute \(s=t_i-t_j\ge0\), collapse the two coordinates, and extend the gap integral to infinity. The Jacobian in \((t_i,s)\) is one; the error in the remaining Lipschitz maxima is bounded by \(|s|\) times a polynomial of one lower degree, and the endpoint error is exponentially anchored. Thus the leading multiplier for a collapsed tie is exactly \(2w_0\), with all these errors \(O(T^{n+2})\) on the fifth moment.

A universal tie collides two global maximizers. The sum of the two-upper-face polynomials equals half of
\(\partial_t^2(t^n Z_J(zt))\big|_{t=1}
+\frac{dz}{1-z}\partial_t(t^n Z_J(zt))\big|_{t=1}\):
on a single upper face precisely \(d\) incident factors equal \(1-zt\) and give the stated derivative correction. Hence its coefficient of degree five again cancels exactly. The only leading soft term comes from an old marked edge \(e\) whose endpoints collide. Delete \(e\) and merge its endpoints; every other edge survives with its multiplicity. The degree-four signed collapsed integral is therefore
\[
 C_{\rm mark}=[z^4]\Delta\sum_e Z_{J/e}(z).
\]
The same recurrence above with multiplicity gives \(1/4725\) for the D contrast and \(-1/1890\) for the C contrast; all lower marked differences vanish by lower shape equality. With the unmarked lower coefficients matched, \(C_{\rm mark}=5/4725-1/1890=1/1890\). The bundled exact replay computes these rational coefficients using integer multiples of \((|S|+q)!\).

The collapsed dimension is \(n-1\) and the polynomial degree four, giving radial factor \((n+3)!\). Replacing the common radial weight by \(e^{-t/T}\) inside these surviving scaled integrals is now justified by dominated convergence. The five reference factors and the common site density cancel the powers of \(h_0\); the remaining denominator is \(n!\). This gives \(\mathcal B=2w_0(n+3)!C_{\rm mark}/n!\). The fifth moment residual is \(o(T^{n+3})\) after the radial limit, the sixth has its dominated hard limit, and all comparisons used the signed finite subtraction. Both scalar leading terms are positive by Proposition~\ref{prop:thermal}, so these are actual relative moduli. Finally \(g^{-5}T^{-2}/g^{-6}=g^3/L^2\), proving the sign change.
\end{proof}

\section{Signed rigid coefficients under dense perturbations}
\label{den:section}

The logarithmic comparisons above detect differences within growing families.
We next consider a different question: whether the signed leading coefficient of a rigid core is stable under changes to a dense graph. The two limits in this section are ordered. First take the WRT level to infinity for the fixed graph, as in Section~\ref{sec:sewing}. Then let the graph size tend to infinity. On a connected nonbipartite core with positive fractional edge weights of incident sum one, the rigid signed coefficient, apart from the prefactor fixed by the vertex and edge counts, is
\[
 J(G)=\int_{(0,\infty)^{V(G)}}\prod_{ij\in E(G)}s(x_i x_j)\prod_i dx_i,
 \qquad s(z)=\frac{\sin(\pi z)}{\pi z}.
\]
The oscillatory factors in this integral require control of the exterior in absolute value, even though the eventual sign below is positive.

\begin{theorem}[Vertex saddle for sublinear deletion degree]\label{den:thm}
Let \(H\) be a fixed connected nonbipartite simple graph, and form a blowup \(G_n^0\) by replacing class \(a\) by \(n_a\) independent vertices, with complete connections on the edges of \(H\). Assume \(n=\sum_a n_a\), \(n_a/n\to\alpha_a>0\), and that there are strictly positive numbers \(w_{ab}=w_{ba}\) on \(E(H)\) with \(\sum_{b:ab\in E(H)}w_{ab}=\alpha_a/2\). Delete any set \(F_n\) of edges of maximum degree \(\Delta_n=o(n)\), obtaining \(G_n\). Put \(\beta=\pi^2/6\). For all sufficiently large \(n\) there is a unique positive solution of
\begin{equation}\label{den:site}
 \frac{\beta}{n}m_i\sum_{j:ij\in E(G_n)}m_j=\frac12 .
\end{equation}
With \(a_i=(2m_i)^{-1}\), define
\[
 {\cal A}_n=n^{-n/4} e^{n/4}\prod_i\frac{\sqrt\pi}{2\sqrt{a_i}}.
\]
Let \(t_a>0\) be the unique solution of
\(\beta t_a\sum_{b:ab\in E(H)}t_b=\alpha_a/2\), let \(A_H\) be the ordinary adjacency matrix, and set \(\Sigma=\operatorname{diag}(2t_a^2/\alpha_a)\). Then
\begin{equation}\label{den:result}
 \frac{J(G_n)}{{\cal A}_n}\longrightarrow
 \frac{\displaystyle
   \exp\!\left[-\frac{\pi^4}{20}\sum_{ab\in E(H)}
                             \frac{t_a^2t_b^2}{\alpha_a\alpha_b}\right]}
 {\sqrt{\det(I+\beta\Sigma^{1/2}A_H\Sigma^{1/2})}}>0 .
\end{equation}
The matrix in the denominator is positive definite. No pairwise degree balance, grouping of removed edges, or rate condition beyond \(\Delta_n=o(n)\) is needed.
\end{theorem}

\subsection{The template saddle and uniform residual comparisons}

Write \(B\) for the signless incidence matrix of \(H\), with column \(e_a+e_b\) at \(ab\). It has full row rank: an element of \(\ker B^t\) alternates sign across every edge, so an odd cycle and connectedness force it to vanish. Using a fixed right inverse of \(B\), the positive weights in Theorem~\ref{den:thm} extend uniformly to all fractions \(\gamma\) near \(\alpha\), with \(Bw=\gamma/2\). The log-radius exponent
\[
 f_\gamma(u)=\tfrac12\sum_a\gamma_a u_a-\beta\sum_{ab\in E(H)}e^{u_a+u_b}
\]
is a sum over edges of \(w_{ab}(u_a+u_b)-\beta e^{u_a+u_b}\). The weights are uniformly positive and the pair sums span all linear directions. Consequently \(f_\gamma(u)\le C-c\sum_a|u_a|\), and its Hessian is strictly negative definite. Its unique uniformly bounded maximum has \(t_a=e^{u_a}\) solving the corresponding saddle equations. These assertions also hold with \(\beta/2\) in place of \(\beta\).

We will restrict to residual graphs obtained by removing \(\ell=o(n)\) vertices, keeping the scale \(n\) in every equation. Write \(r_a\) for residual class sizes, \(N'=\sum r_a\), and \(\gamma_a=r_a/n\). For the bare residual template, the unique site solution is constant in each class, \(m_a^0=e^{u_a}/\gamma_a\). In logarithmic site coordinates the relevant functional is
\[
 {\cal F}_G(y)=\tfrac12\sum_i y_i-\frac{\beta}{n}
                  \sum_{ij\in E(G)}e^{y_i+y_j}.
\]
Its negative Hessian at the bare solution is the weighted signless Laplacian. On each class-mean-zero subspace it is multiplication by \(1/2\); on vectors constant within classes it is uniformly invertible by the same signless rank condition, with all class fractions bounded away from zero. The projections onto these subspaces are bounded in sup norm, so the bare inverse is uniformly bounded in that norm. Deletion changes both the gradient and the Hessian at the bare point by \(O(\Delta_n/n)\) in the respective sup norms, and the Hessian has uniformly bounded local Lipschitz constant there by its row-sum formula. The contraction obtained by subtracting the bare gradient linearization yields a stationary point with
\begin{equation}\label{den:vertexclose}
 \max_i|m_i-m_{a(i)}^0|=O(\Delta_n/n).
\end{equation}
For example, fixing the bare inverse, the Newton correction map sends a sup-norm ball of radius \(C\Delta_n/n\) into itself and has Lipschitz constant \(O(\Delta_n/n)\) on this ball. If there are no deletions the bare solution suffices.

The active graph retains signless coercivity
\begin{equation}\label{den:coercivity}
 \sum_{ij\in E(G)}|y_i+y_j|\ge c n\sum_i |y_i|.
\end{equation}
For the bare graph, let \(\bar y_a\) be the class mean and \(D_{ab}=\sum_{i\in a,j\in b}|y_i+y_j|\) on a template edge. Jensen gives \(r_ar_b|\bar y_a+\bar y_b|\le D_{ab}\) and
\(\sum_{i\in a}|y_i-\bar y_a|\le2D_{ab}/r_b\).
Sum the latter over classes using one adjacent class each, and control the class means by injectivity of \(B^t\). This proves \eqref{den:coercivity} before deletion; the deletion loss is at most \(\Delta_n\sum_i|y_i|\). Thus the Hessian of \({\cal F}_G\) is strictly negative everywhere, making the nearby point its unique global maximum. Connectivity also persists: within any class two vertices retain a common neighbor in an adjacent class, and adjacent classes retain edges. Equation \eqref{den:site} gives actual positive fractional weights \(p_{ij}=(2\beta/n)m_i m_j\asymp n^{-1}\) on all active edges, of incident sum one.

For a residual graph define \({\cal A}_G\) by replacing the site count \(n\) by \(N'\) in the powers and the product of \({\cal A}_n\), retaining the scale \(n^{-N'/4}\).
Summing \eqref{den:site} gives \((\beta/n)\sum_{ij\in E(G)}m_i m_j=N'/4\).
Hence for a fixed numerical \(c_0\),
\[
 \log{\cal A}_G=\max{\cal F}_G+N'/2-(N'/4)\log n+c_0N'.
\]
All optimizing \(m_i\) are uniformly bounded above and away from zero. Comparing the maximum before and after restriction by evaluating each candidate on the other's domain, extending by fixed bounded coordinates in the reverse direction, changes it by \(O(\ell)\). Each removed site affects only one log term and \(O(n)\) bounded edge terms divided by \(n\). It follows that
\begin{equation}\label{den:resratio}
 {\cal A}_{G}/{\cal A}_n\le(Cn^{1/4})^\ell,\qquad
 {\cal A}_n\ge n^{-n/4}e^{-Cn}.
\end{equation}
Removing edges increases \({\cal F}_G\) pointwise, so \({\cal A}_G\) is at least the corresponding bare residual normalizer.

We also use two bare radius bounds. If \(T_a=\sum_{i\in a}x_i^2\) and \(V^0=\beta\sum_{ab\in E(H)}T_aT_b\), polar coordinates by classes put the radius integral for \(Q^0=\int e^{-V^0}dx\) in the form \(\int e^{n f_\gamma(u)}du\), up to the radius-independent prefactor, with \(u_a=\log(T_a/\sqrt n)\). The uniform linear bound above implies that outside a sufficiently wide fixed box \(a_0\le T_a/\sqrt n\le A_0\) the radius integral at coefficient \(\beta/2\) can be bounded by \(e^{-Ln}\) times its full integral at coefficient \(\beta\), for any prescribed \(L\). Indeed reserve half the negative linear bound for integration, and lower bound the denominator on a fixed compact box. Below we also obtain \(Q^0\le C{\cal A}_G\).

\subsection{Exact Gamma tilt and the compact fluctuation estimate}

Put \(U_i=\sqrt n\,x_i^2\). Under \(P_m\) let the \(U_i\) be independent Gamma variables with shape \(1/2\) and rate \(a_i\), thus mean \(m_i\). Define
\[
 X_a=n^{-1/2}\sum_{i\in a}(U_i-m_i),\quad
 B_0(X)=\beta\sum_{ab\in E(H)}X_aX_b,\quad
 L_F=\frac{\beta}{n}\sum_{ij\in F}(U_i-m_i)(U_j-m_j).
\]
In all formulas on residual sets, \(F\) is the restricted deletion graph. The exact change of measure for the surviving quartic energy is
\begin{equation}\label{den:tilt}
 e^{-(\beta/n)\sum_{ij\in E(G)}U_iU_j}\prod_i dx_i
       ={\cal A}_G e^{-B_0(X)+L_F}\,dP_m .
\end{equation}
The constant term of the energy at \(m\) is \(N'/4\); its linear term is \(\sum_i a_iU_i-N'/2\) by \eqref{den:site}; its centered quadratic is \(B_0-L_F\). The Gamma factors in \eqref{den:tilt} are exactly those in \(dx_i=n^{-1/4}dU_i/(2\sqrt{U_i})\).

For a fixed compact radius box enclosing the limiting saddle there is \(p_0>1\) such that, uniformly for the residual sets under consideration,
\begin{equation}\label{den:block}
 \mathbb E_m[e^{-p_0 B_0(X)}\mathbf 1_{\rm box}]\le C .
\end{equation}
To prove this even when \(B_0\) is indefinite, write \(v_a=T_a/\sqrt n=n^{-1}\sum_{i\in a}U_i\) and \(v^*_{a,n}=n^{-1}\sum_{i\in a}m_i\). The class log-mgf divided by \(n\) is \(\psi_{a,n}(s)=-(2n)^{-1}\sum_{i\in a}\log(1-2m_i s)\). On the box, the one-dimensional tilt \(s\) satisfying \(\psi'_{a,n}(s)=v_a\) stays in a compact interior interval of its domain. This follows from \eqref{den:vertexclose}: all means in class \(a\) tend uniformly to \(\mu_a=t_a/\alpha_a\), and the limiting derivative \(t_a/(1-2\mu_a s)\) ranges over the positive half-line. Under the tilted law, the density of the scaled class sum \(X_a\) is bounded by a uniform constant. Fourier inversion gives this since the modulus of its characteristic function is at most \((1+c\xi^2/n)^{-r_a/4}\); its integral in \(\xi\) is \(O(1)\). Undoing the tilt bounds the density at \(X_a=\sqrt n(v_a-v_{a,n}^*)\) by \(C\exp[-n I_{a,n}(v_a)]\), where \(I_{a,n}=sv_a-\psi_{a,n}(s)\). The rates converge in \(C^2\) on the box to
\[
 I_a(v)=\frac{v-t_a}{2\mu_a}-\frac{\alpha_a}{2}\log(v/t_a).
\]
They vanish to first order at their finite centers. At \(p=1\), the limit of
\(\sum_a I_{a,n}(v_a)+p\beta\sum_{ab}(v_a-v_{a,n}^*)(v_b-v_{b,n}^*)\)
is the strict log-radius objective gap,
\[
 \beta\sum_{ab}(v_av_b-t_at_b)-\tfrac12\sum_a\alpha_a\log(v_a/t_a).
\]
The identity uses \(1/(2\mu_a)=\beta\sum_{b:ab\in E(H)}t_b\). Its Hessian at \(v=t\) is \(\Sigma^{-1}+\beta A_H>0\) by the strictly negative log-radius Hessian. A compact annular gap, uniform \(C^2\) convergence and vanishing first derivatives at the finite centers give a lower bound \(c\|v-v_n^*\|^2\) on the box also for some \(p_0>1\). Apply the independent density bounds and integrate \(e^{-c\sum X_a^2}\) to obtain \eqref{den:block}. This also proves the positive definiteness in the theorem. Applied to the bare graph, \eqref{den:tilt} and \eqref{den:block} give \(Q^0\le C {\cal A}_{G^0}\): the compact box contains at least half the bare radius mass by its Laplace estimate. Thus \(Q^0\le C{\cal A}_G\) as asserted above.

\subsection{Neighborhood control of the deleted quadratic}

Fix a small constant \(\eta>0\) and use the event
\begin{equation}\label{den:admissible}
 {\cal E}_G=\left\{\sum_{j:ij\in F}U_j\le\eta n
                         \text{ for every residual site }i\right\}.
\end{equation}
For each fixed \(p>0\), with \(\eta\) sufficiently small,
\begin{equation}\label{den:mgf}
 \mathbb E_m[e^{p|L_F|}\mathbf 1_{{\cal E}_G}]\le C_p .
\end{equation}
The site Gamma laws have the uniform centered estimate
\(\log\mathbb E e^{s(U_i-m_i)}\le C s^2\) for \(|s|\le s_0\), by their bounded means and rates. Consider a single deleted class-pair matrix \(A\) and the independent centered Gamma vectors \(\xi,\zeta\) at its two endpoints. Retaining only the restrictions \({\cal E}_a=\{\sum_i A_{ij} U_i\le\eta n\ \forall j\}\) from the first side, condition on \(\xi\) to bound
\[
 \mathbb E[e^{(\theta/n)\xi^t A\zeta}\mathbf 1_{{\cal E}_G}]
 \le \mathbb E_\xi[e^{C_\theta\|A^t\xi\|^2/n^2}\mathbf 1_{{\cal E}_a}]
\]
for fixed \(\theta\in\mathbb R\). Indeed \(\|A^t\xi\|_\infty\le\eta n+C\Delta_n\), within the local mgf interval after division by \(n\) and a suitable choice of \(\eta\).

Linearize the square using a standard Gaussian \(g\) on the column-label space. Require \(\|Ag\|_\infty\le c_1n\) for a small fixed \(c_1>0\). Completing the Gaussian square at each \(\xi\in{\cal E}_a\) shifts \(g\) by \(g_0=\sqrt{2C_\theta}A^t\xi/n\). Then \(\|Ag_0\|_\infty=O(\Delta_n(\eta+\Delta_n/n))=o(n)\), and each unshifted coordinate of \(Ag\) has variance at most \(\Delta_n\). The shifted restriction therefore retains \(1-o(1)\) of its Gaussian mass, uniformly in \(\xi\), by a union bound \(O(n e^{-c n^2/\Delta_n})\); the empty matrix is immediate. Reverse the two expectations after this restriction, dropping the \(\xi\)-event and now using the local Gamma mgf on the first side. The right-hand side is bounded by
\[
 (1-o(1))^{-1}\mathbb E_g e^{C'_\theta\|Ag\|^2/n^2}=1+o(1).
\]
The Gaussian determinant is bounded because \(\|A\|_{\mathrm{op}}=O(\Delta_n)\) and \(\|A\|_F^2=O(n\Delta_n)\); the associated eigenvalue bound and trace bound tend to zero. Taking both signs and using finite H\"older across class pairs proves \eqref{den:mgf}. The argument does not assume a finite unconditional exponential moment for a product of uncapped Gamma variables.

Under the original product Gamma law, \(P_m({\cal E}_G^c)=o(1)\): a fixed positive-parameter Chernoff estimate for each missing-neighbor sum is \(e^{-c\eta n+C\Delta_n}\). Also \(L_F\to0\) in probability by its variance \(O(\Delta_n/n)\), since centered deleted-edge products sharing only one endpoint have zero covariance. From \eqref{den:block} and \eqref{den:mgf} we obtain a uniform \(L^p\) bound, for some \(p>1\), on \(e^{-B_0+L_F}\) restricted to \({\cal E}_G\) and the compact radius box. Choose the first H\"older exponent sufficiently close to one and \(p\) small enough so their product is less than \(p_0\), then fix \(\eta\) to accommodate the other exponent. The radius box can be fixed before \(\eta\), as made explicit in the tail bounds next.

\subsection{Absolute exterior at every sublinear deletion rate}

Use \(x_i=n^{-1/4}e^{y_i}\), so \(U_i=e^{2y_i}\), and the global bound \(|s(z)|\le (1+2\beta z^2)^{-1/2}\). For \(z\le1\) this follows from Euler's product and \(\log(1+u)\le u\); for \(z\ge1\) use \(|s(z)|\le1/(\pi z)\). The log-envelope potential satisfies, uniformly for \(p_{ij}\asymp n^{-1}\),
\[
 \tfrac12\log(1+2\beta e^{2v}/n)-p_{ij}v
 \ge \frac d n |v|-\frac C n+
       \frac {c_3}n\mathbf1_{v\ge0}\min(e^{2v},n).
\]
On \(v<0\) use nonnegativity and the lower bound on \(np_{ij}\). Up to \(v=(\log n)/2\) use the exponential lower bound, retaining a fixed fraction of it and absorbing the desired linear term in the rest. Above that point use the positive derivative, bounded below independently of \(n\). Summing, the fractional linear term cancels the Jacobian, and \eqref{den:coercivity} gives the pointwise absolute bound
\begin{equation}\label{den:log-envelope}
 n^{-n/4}\exp\!\left[
   Cn-c_4\sum_i|y_i|-\frac{c_5}{n}
     \sum_{ij\in E(G_n)}
       \mathbf1_{y_i+y_j\ge0}\min(e^{2(y_i+y_j)},n)\right]\,dy .
\end{equation}
This already proves absolute convergence at each sufficiently large \(n\).

Let \(c_6n\) be a uniform lower bound on all active degrees. The part with at least \(c_6n/4\) sites \(y_i<-A\) has absolute mass at most \(e^{-n}{\cal A}_n\) for sufficiently large fixed \(A\): extract half the independent log decay in \eqref{den:log-envelope} before integrating the rest and use \eqref{den:resratio}. On the remaining configurations every site has a positive fraction of \(n\) retained neighbors with \(U_j\ge e^{-2A}\). If \(M\to\infty\), \(M\le n\), and at least \(D n/M\) sites satisfy \(U_i>M\), the last term of \eqref{den:log-envelope} gives another \(e^{-n}{\cal A}_n\) bound for large fixed \(D\): charge \(cM\) to each such site using its non-tiny neighbors, counting each edge at most twice. The constants are independent of the rate at which \(M\) grows.

Otherwise fix the high set above \(M\), of size \(\ell<Dn/M=o(n)\). Each high site still has a positive fraction of \(n\) retained \emph{low} neighbors with \(U_j\ge e^{-2A}\), since the removed high set is small. Its high-low factors are bounded by \((1+\kappa U_i/n)^{-\nu n}\) for fixed \(\kappa,\nu>0\). These edge sets are disjoint for different high sites. Half of the bound gives \(e^{-c\min(U_i,n)}\), and the other half has bounded integral in \(dU/(2\sqrt U)\), since
\[
 \int_0^\infty (1+\kappa U/n)^{-\nu n/2}\frac{dU}{2\sqrt U}
  =\tfrac12\sqrt{n/\kappa}\,
    \mathrm B(1/2,\nu n/2-1/2)\le C.
\]
Thus on a set forcing \(\sum_{\rm high}\min(U_i,n)\ge R_*\), integrate the high variables at cost \(e^{-cR_*}(Cn^{-1/4})^\ell\). Extract the penalty uniformly before dropping constraints on high variables; no hypothesis of independent constraints is needed. The \(n^{-1/4}\) is the physical \(x\)-measure factor, and cancels the factor \(n^{1/4}\) per removed site in \eqref{den:resratio}.

First bound \({\cal E}_{G_n}^c\) by taking
\(M=\sqrt{n/\max(\Delta_n,\sqrt n)}\).
Then \(\Delta_n M=o(n)\), so a violating neighborhood has at least \(\eta n/2\) of high charge; the aggregate \(\sum_{\rm high}\min(U_i,n)\) is at least \(\eta n/2\) if \(\eta\le1\). The number of high sets, including \(C^\ell\), is \(\exp(o(n))\) by the binomial entropy bound since \(\ell<Dn/M\). For the low variables write \(S_a=\sum_{i\in a,{\rm low}}U_i\). Their products are in the positive sine regime since \(M=o(\sqrt n)\), so their product is bounded above by \(e^{-V^0+D_{\rm lost}}\). For each pair of classes its lost energy is at most \((\beta/n)\Delta_n M\min(S_a,S_b)\). On the compact radius box this is \(o(n)\). Outside it, the bound \(D_{\rm lost}\le V^0/2+C(\Delta_n M)^2/n=V^0/2+o(n)\) follows by separating \(\max(S_a,S_b)\ge2\Delta_n M\) from its complement for each pair. The bare radius estimates and \(Q^0\le C{\cal A}_{G}\) now apply. On the box the high penalty beats \(e^{o(n)}\) times the residual normalization; outside it, the broad bare radius box alone gives exponential suppression. Thus the absolute integral on \({\cal E}_{G_n}^c\) is \(o({\cal A}_n)\).

On \({\cal E}_{G_n}\) take instead a fixed positive cube \(x_i\le b\), \(b^2\le1/2\), with high threshold \(M=b^2\sqrt n\). At least \(D n/M\) high sites were already excluded by \eqref{den:log-envelope}. For fewer than this and at least one high site, integrate the high variables as above at cost \((Cn^{-1/4}e^{-c\sqrt n})^\ell\). The low-site neighborhood bound \({\cal E}_G\) is inherited by removal of high sites. On the radius box, \eqref{den:tilt} and its \(L^p\) bound control the low product by \(C{\cal A}_G\) after integration. Summing high sets costs at most \(\sum_{\ell\ge1}(Cn e^{-c\sqrt n})^\ell=o(1)\) relative to \({\cal A}_n\). Outside the box, use the neighborhood bound to get \(D_{\rm lost}\le V^0/2+C\eta^2n\) by the same pairwise argument with \(\eta n\) in place of \(\Delta_n M\). The bare box can be chosen first large enough to beat \(C\eta^2 n\) for every \(\eta\le1\). Its tail estimate then applies also at \(\ell=0\). This proves that everything outside the intersection of \({\cal E}_{G_n}\), the fixed cube, and the compact radius box contributes \(o({\cal A}_n)\) in absolute measure.

\subsection{Evaluation of the signed integral}

Under \(P_m\), the fixed cube has probability tending to one by the union bound for Gamma tails, the neighborhood restriction was checked above, and the compact box has probability tending to one by concentration of class sums. Equation \eqref{den:vertexclose} and the independent triangular-array central limit theorem give \(X\Rightarrow N(0,\Sigma)\), with diagonal covariance as in Theorem~\ref{den:thm}. On the cube the product divided by its surviving quartic envelope is positive and at most one, and its logarithm is
\[
 -\frac{\pi^4}{180}\sum_{ij\in E(G_n)}x_i^4x_j^4
       +O\big((\sum_i x_i^6)^2\big).
\]
Under the product Gamma law \(\sum_{i\in a}x_i^4\to3t_a^2/\alpha_a\) by the law of large numbers, \(\sum_i x_i^6=O_P(n^{-1/2})\), and \(\sum_i x_i^8=O_P(n^{-1})\). The deleted fourth-order terms are at most \(C\Delta_n\sum_i x_i^8=o_P(1)\), so the sine correction tends to the exponential in \eqref{den:result}. Combine the block convergence, \(L_F\to0\), and the restricted \(L^p\) bound to take expectations in \eqref{den:tilt}. The Gaussian expectation of \(e^{-B_0}\) is the inverse square root of \(\det(I+\beta\Sigma^{1/2}A_H\Sigma^{1/2})\). The absolute exterior bound completes the proof of Theorem~\ref{den:thm}.

\subsection{Linear-degree deletion: a loop permutation seen by the determinant}

We give a boundary example in which vertex degrees, pair densities, the product saddle normalization, and the leading WRT phase can all be held fixed. Start with three coarse classes of size \(qN\) each, equip each with \(q\) microclasses of \(N\) vertices, and fix \(q\ge3\). Across each coarse pair remove the \(K_{N,N}\)'s specified by a matching of the \(q\) labels. In each class the labels may be permuted independently. Thus the first two removed matchings can be taken as the identity and the remaining one is a permutation \(P\), the ordered product of the three original matchings around the triangle (up to conjugation, or inversion upon reversing the triangle). This is a statement about the retained adjacency, not a continuum gauge interpretation.

\begin{proposition}[Loop permutation determinant]\label{den:loop}
For the resulting graphs \(G_{N,P}\), all sites lose \(2N\) neighbors and the deleted fraction between each coarse pair is \(1/q\). Write \(k_0=3q\), \(d_0=2(q-1)\), and \(T_P\) for the retained quotient on microclasses. The product normalization \({\cal A}_n\) is independent of \(P\), and
\begin{equation}\label{den:det-loop}
 \begin{aligned}
 \frac{J(G_{N,P})}{{\cal A}_n}
 &\longrightarrow e^{-9k_0/(40d_0)}\det(I+A_{T_P}/d_0)^{-1/2},\\
 \det(I+A_{T_P}/d_0)
 &=\frac12\prod_{e^{i\theta}\in\operatorname{spec}(P|_{\mathbf1^\perp})}
              \left(1-\frac3{d_0^2}-\frac{2\cos\theta}{d_0^3}\right).
 \end{aligned}
\end{equation}
The product is over eigenvalue multiplicity.
\end{proposition}

\begin{proof}
The refined quotient is \(d_0\)-regular on \(k_0\) vertices. Across a pair of coarse classes the quotient is \(K_{q,q}\) minus a matching, which is connected for \(q\ge3\) since any two vertices on one side have a common neighbor. An allowed \(A\)--\(C\) pair of labels, together with a \(B\) label different from both, gives a triangle. Thus \(T_P\) is connected and nonbipartite, with positive fractional weights \(1/d_0\). Theorem~\ref{den:thm} applies to its \emph{bare} equal-size blowup, rather than to the linear deletions from the original coarse graph. The radial solution is \(t_a^2=(2\beta k_0 d_0)^{-1}\); the common site mean is \(\sqrt{k_0/(2\beta d_0)}\); the covariance is \(\Sigma=(\beta d_0)^{-1}I\). This gives the first formula of \eqref{den:det-loop}, including its \(P\)-independent sine correction.

The three between-class adjacency blocks are \(J_q-I_q,J_q-I_q,J_q-P\), up to transposes; here \(J_q\) is the all-ones matrix. On the uniform label mode the normalized three-by-three block \(I+A_{T_P}/d_0\) has diagonal one and off-diagonal \(1/2\), giving determinant \(1/2\). On the orthogonal label space \(J_q=0\). Diagonalizing \(P\) over \(\mathbb C\) gives for each eigenvalue \(e^{i\theta}\) a three-by-three Hermitian block with off-diagonals \(-1/d_0,-1/d_0,-e^{i\theta}/d_0\) in a cyclic orientation, up to conjugation of the last entry. The determinant expansion gives the second formula.
\end{proof}

The cycle dependence can also be written in real terms. Put \(a=1-3/d_0^2\), \(b=2/d_0^3\), \(a_*=(a+\sqrt{a^2-b^2})/2\), and \(v=b/(2a_*)\in(0,1)\). Then \(a-b\cos\theta=a_*|1-v e^{i\theta}|^2\). For cycle lengths \(\ell_1,\ldots,\ell_j\) of \(P\), taking products over the roots of unity and removing one uniform mode gives
\[
 \det(I+A_{T_P}/d_0)=
        \tfrac12 a_*^{q-1}\frac{\prod_i(1-v^{\ell_i})^2}{(1-v)^2}.
\]
Merging two cycles strictly increases the determinant since
\((1-v^u)(1-v^w)<1-v^{u+w}\) for positive \(u,w\), and decreases the signed coefficient's normalized limit. This dependence persists at arbitrarily small fixed positive deleted fraction \(2/(3q)\).

For \(q=4\), compare \(P\) of type \((2,2)\) with a four-cycle. The normalized signed leading ratio from \eqref{den:det-loop} is
\[
 \frac{a}{\sqrt{a^2-b^2}}=\frac{99}{70\sqrt2}>1 .
\]
The quotient adjacency signatures in both cases are \(-2\). Indeed the uniform block has signature \(-1\); a nonuniform adjacency block has trace zero and determinant \(-2\cos\theta\), so its signature is respectively \(+1,-1,0\) for positive, negative, or zero \(\cos\theta\). The extra zero modes of the blowup do not change signature. The two graphs also have exactly the same edge and vertex counts. Thus the ratio distinguishes even the coefficients with the same oriented WRT phase under the convention of Section~\ref{sec:sewing}. It is not an effect of the Maslov factor. Alternating the two refinements yields no common vertex-normalized limit at linear deletion degree despite identical degree, density, site-mean, and phase data.

\paragraph{Order of limits.}
The vertex theorem does not interchange the WRT and graph-size limits. The finite expression already shows why uniformity would need separate estimates. For example, in the balanced tripartite graph \(K_{m,m,m}\), keeping the original odd charge equal to one on two entire classes and summing the third freely produces a nonnegative sector of the normalized finite sum of size
\((1+o(1))\,2^{n/3}r^{-n/6}\) when \(n=3m\) and \(r\asymp n^{3/2}\). For this sector the only edge factor after cancellation is \(R(1/r)= (\pi/r)/\sin(\pi/r)\) per edge, and each free site has cap sum \(\sum_{k<r,\ k\ {\rm odd}}\cos^2(\pi k/(2r))=r/4\) on either parity. The sector has the same \(n^{-n/4}\) dimensional power as the tripartite continuum bulk on that scale. It is not a lower bound on the full signed sum.

\section{What survives}
Equal low log coefficients need not erase an incidence motif. Joint charge weighting changes the size of a graph contrast and, for a suitable pair, reverses its sign. In dense cores arbitrary \(o(n)\) maximum-degree edge deletion becomes invisible in the signed fluctuation multiplier after vertexwise normalization, not in the raw invariant. At linear deletion, loop ordering changes the normalized determinant. This determinant is not an APS invariant or a twist of a marked operator. The conclusions are neither reconstruction at all orders nor a general graphon limit, and the sequential density theorem is not a simultaneous level theorem.

\begin{thebibliography}{9}
\bibitem{BHMV} C. Blanchet, N. Habegger, G. Masbaum and P. Vogel, Topological quantum field theories derived from the Kauffman bracket, \emph{Topology} 34 (1995), 883--927.
\bibitem{Turaev} V. G. Turaev, \emph{Quantum Invariants of Knots and 3-Manifolds}, De Gruyter, 1994, Chapter IV.
\bibitem{GM} P. M. Gilmer and G. Masbaum, Maslov index, Lagrangians, mapping class groups and TQFT, \emph{Forum Math.} 25 (2013), 1067--1106; \href{https://arxiv.org/abs/0912.4706}{arXiv:0912.4706}.
\end{thebibliography}
\end{document}
